\begin{solutionbox}
\textbf{Given.} The singlet state $\ket{\psi}=\frac{1}{\sqrt2}\big(\ket{\uparrow\downarrow}-\ket{\downarrow\uparrow}\big)$, written $\ket{\text{Alice}\,\text{Bob}}$. Alice's operators $\sigma_i$ act only on the first label and Bob's $\tau_i$ only on the second (formally $\sigma_x\otimes\mathbb{I}$ and $\mathbb{I}\otimes\tau_y$). As in the Week 7 notes they are Pauli operators with eigenvalues $\pm1$:
\[
\sigma_x\ket{\uparrow}=\ket{\downarrow},\quad\sigma_x\ket{\downarrow}=\ket{\uparrow};\qquad
\tau_y\ket{\uparrow}=i\ket{\downarrow},\quad\tau_y\ket{\downarrow}=-i\ket{\uparrow}.
\]

\textbf{Order of measurement.} The question specifies that Bob measures first. Because $\sigma_x$ and $\tau_y$ act on different subsystems they commute, $[\sigma_x\otimes\mathbb{I},\mathbb{I}\otimes\tau_y]=0$, so the joint statistics, and hence $\langle\sigma_x\tau_y\rangle$, are the same in either order. The order is therefore irrelevant for every quantity asked for.

\textbf{(a) $\boldsymbol{\langle\sigma_x\tau_y\rangle}$.} Apply $\tau_y$ (Bob) first, then $\sigma_x$ (Alice), term by term:
\begin{align*}
\tau_y\ket{\uparrow\downarrow}=-i\ket{\uparrow\uparrow}&\ \Rightarrow\ \sigma_x\tau_y\ket{\uparrow\downarrow}=-i\ket{\downarrow\uparrow},\\
\tau_y\ket{\downarrow\uparrow}=i\ket{\downarrow\downarrow}&\ \Rightarrow\ \sigma_x\tau_y\ket{\downarrow\uparrow}=i\ket{\uparrow\downarrow}.
\end{align*}
Hence
\[
\sigma_x\tau_y\ket{\psi}=\frac{1}{\sqrt2}\big(-i\ket{\downarrow\uparrow}-i\ket{\uparrow\downarrow}\big)=-\frac{i}{\sqrt2}\big(\ket{\uparrow\downarrow}+\ket{\downarrow\uparrow}\big).
\]
Taking the inner product with $\bra{\psi}=\frac{1}{\sqrt2}\big(\bra{\uparrow\downarrow}-\bra{\downarrow\uparrow}\big)$ and using orthonormality of the four basis states:
\[
\langle\sigma_x\tau_y\rangle=-\frac{i}{2}\big(\braket{\uparrow\downarrow|\uparrow\downarrow}-\braket{\downarrow\uparrow|\downarrow\uparrow}\big)=-\frac{i}{2}(1-1)=0 .
\]
The operator maps the singlet onto the (orthogonal) triplet state $\frac{1}{\sqrt2}(\ket{\uparrow\downarrow}+\ket{\downarrow\uparrow})$, so the overlap vanishes.

\textbf{(b) Single-spin expectation values.}
\begin{align*}
\tau_y\ket{\psi}&=\tfrac{1}{\sqrt2}\big(-i\ket{\uparrow\uparrow}-i\ket{\downarrow\downarrow}\big)\ \Rightarrow\ \langle\tau_y\rangle=0,\\
\sigma_x\ket{\psi}&=\tfrac{1}{\sqrt2}\big(\ket{\downarrow\downarrow}-\ket{\uparrow\uparrow}\big)\ \Rightarrow\ \langle\sigma_x\rangle=0,
\end{align*}
since in each case the result contains only $\ket{\uparrow\uparrow}$ and $\ket{\downarrow\downarrow}$, which are orthogonal to both terms of $\ket{\psi}$.

\textbf{(c) Correlation.}
\[
\langle\sigma_x\tau_y\rangle-\langle\sigma_x\rangle\langle\tau_y\rangle=0-0\cdot0=0 .
\]

\textbf{Second method (sequential collapse).} Bob's $\tau_y$ outcome is $\pm1$ with probability $\tfrac12$ each (since $\langle\tau_y\rangle=0$). If Bob obtains $\tau_y=+1$, the singlet's perfect anticorrelation along any common axis means Alice's spin collapses to the $\sigma_y=-1$ eigenstate $\frac{1}{\sqrt2}(\ket{\uparrow}-i\ket{\downarrow})$, for which $\langle\sigma_x\rangle=0$; likewise for $\tau_y=-1$. Alice's $x$-outcome is therefore $\pm1$ with equal probability whatever Bob found, so $\langle\sigma_x\tau_y\rangle=\sum P(s_A,s_B)\,s_As_B=0$, confirming (a)--(c).

\textbf{Checks.} With $4\times4$ Kronecker-product matrices: $\langle\sigma_x\tau_y\rangle=\langle\sigma_x\rangle=\langle\tau_y\rangle=0$, the two operators commute, the full correlation matrix is $\langle\sigma_i\tau_j\rangle=-\delta_{ij}$, and the projective simulation of ``Bob first, then Alice'' gives zero joint correlation and $\langle\sigma_y\rangle=\mp1$ after Bob's $\pm1$ (perfect anticorrelation on the shared axis, as a control). If spin operators $S=\frac\hbar2\sigma$ were used instead, every result would be scaled by $\hbar/2$ or $\hbar^2/4$ and would still be zero.
\end{solutionbox}

\begin{answerbox}
\[
\text{(a) }\langle\sigma_x\tau_y\rangle=0,\qquad\text{(b) }\langle\tau_y\rangle=0,\ \ \langle\sigma_x\rangle=0,\qquad\text{(c) }\langle\sigma_x\tau_y\rangle-\langle\sigma_x\rangle\langle\tau_y\rangle=0 .
\]
Therefore, in the singlet state Alice's $x$-measurement and Bob's $y$-measurement each give $\pm1$ at random and are completely uncorrelated, with every expectation value and the correlation equal to zero, regardless of who measures first.
\end{answerbox}

\begin{insightbox}
Zero correlation here does not mean the spins are not entangled. The singlet is maximally entangled, and its correlations are $\langle\sigma_i\tau_j\rangle=-\delta_{ij}$, i.e.\ $\langle(\hat{\mathbf a}\cdot\boldsymbol\sigma)(\hat{\mathbf b}\cdot\boldsymbol\tau)\rangle=-\hat{\mathbf a}\cdot\hat{\mathbf b}$: perfectly anticorrelated along a common axis, and uncorrelated along perpendicular axes such as $x$ and $y$. Meanwhile each spin on its own is completely random ($\langle\sigma_i\rangle=\langle\tau_i\rangle=0$): we know the joint state exactly yet nothing about either part. That angle-dependent correlation is what Bell tests, and entanglement-based quantum key distribution, rely on.
\end{insightbox}
